
\documentclass[english]{article}
%\documentclass[aps,pra,twocolumn,english]{revtex4}
%\documentclass[english]{iopart}
\usepackage[T1]{fontenc}
\usepackage[latin9]{inputenc}
\usepackage{babel}
\usepackage{graphicx,amsmath,amssymb,color}
\usepackage[normalem]{ulem}
\usepackage{amsfonts}
\usepackage[toc,page]{appendix}
%\usepackage[ocgcolorlinks,colorlinks=true,linkcolor=blue,citecolor=red]{hyperref}
\usepackage{hyperref}
\usepackage{latexsym}
\usepackage{amsfonts}
\usepackage{algpseudocode}
\usepackage{amsthm}
\usepackage{mathrsfs}
\usepackage{natbib}
\usepackage{color,verbatim}
\usepackage{psfrag}
\bibliographystyle{unsrt}
%\bibliographystyle{acm}
%\bibliographystyle{unsrtnat}
%\usepackage[style=numeric]{biblatex}

\usepackage{subcaption} % allows for subfigures


\usepackage[svgnames]{xcolor}
\usepackage{tikz}

\usepackage{algorithm}


%\input{tikzit/tikz_preamble.tex}  % load separately defined tikz preamble


\newcommand{\bra}[1]{\langle#1 |}
\newcommand{\ket}[1]{|#1 \rangle}
\newcommand{\braket}[2]{\left \langle #1 \mid #2 \right \rangle}
\newcommand{\bigbraket}[2]{\left \langle #1 \middle\vert #2 \right \rangle}
\newcommand{\ketbra}[2]{\vert #1 \rangle \! \langle #2 \vert}
\newcommand{\overlap}[2]{\left \langle #1 | #2 \right\rangle}
\newcommand{\average}[1]{\langle #1 \rangle}
\newcommand{\sandwich}[3]{\left \langle #1 \mid #2 \mid #3 \right\rangle}
\newcommand{\innerprod}[2]{\left \langle #1 | #2 \right\rangle}

\begin{document}


\title{Density Operators}
\author{Nicholas Chancellor}


\maketitle


\today % added for drafting, remove in final version


{\small Disclaimer: lecture notes may contain ``typo'' type errors (factors of $2$, minus signs, etc...) if something like this looks wrong there is a good chance it is, please let me know}


\section{How to treat incoherent superpositions of quantum states}

\begin{itemize}
\item Wavevectors don't support incoherent superposition, always a definite phase relationship between different elements (unless one is $0$)
\item Need a representation where probabilities, rather than probability amplitudes sum
\item Outer product of two wavevectors $\ketbra{\psi}{\psi}$, matrix elements can be defined as 
\begin{equation}
\hat{\rho}_{j,k}=\braket{j}{\psi}\braket{\psi}{k}=\psi_j\psi^*_k
\end{equation}
\item Diagonal elements are probabilities $|\psi_k|^2$, superimposing will give sum of probabilities
\item A density operator (sometimes called a density matrix because usually represented as a matrix), can be written as 
\begin{equation}
\hat{\rho}=\sum_j p_j\ketbra{\phi}{\phi}_j
\end{equation}
\item $p_j$ are probabilities of observing different states, note that $\ket{\phi}_j$ \textbf{do not need to be orthogonal}
\item Because density operators are built from copies of state vectors stuck together, it is easy to generalize many operations
\begin{itemize}
\item Operator expectation:
\begin{equation}
\average{\hat{A}}=\sandwich{\psi}{\hat{A}}{\psi}=\mathrm{Tr}[\hat{A}\ketbra{\psi}{\psi}]=\mathrm{Tr}[\hat{A}\hat{\rho}]
\end{equation}
\item $\mathrm{Tr}$ is the trace operation $\mathrm{Tr}[\hat{A}]=\sum_j\sandwich{j}{\hat{A}}{j}$, where $\ket{j}$ is \textbf{any} set states which form a complete orthonormal basis.
\item Schr{\"o}dinger equation\footnote{using theorists units where I have scaled out $\hbar$} $\frac{\partial}{\partial t}\ket{\psi}=-i\hat{H}\ket{\psi}$, apply to left and right spaces separately (note extra minus sign from complex conjugation):
\begin{align}
(\frac{\partial}{\partial t}\ket{\psi})\bra{\psi}+\ket{\psi}(\frac{\partial}{\partial t}\bra{\psi})=-i\hat{H}\ketbra{\psi}{\psi}+i\ketbra{\psi}{\psi}\hat{H} \Rightarrow \\
\frac{\partial}{\partial t}(\ket{\psi}\bra{\psi})=-i[\hat{H},\ketbra{\psi}{\psi}]\Rightarrow \\
i\frac{\partial}{\partial t}\hat{\rho}=[\hat{H},\hat{\rho}]
\end{align}
\item This last equation is known as the von Neumann equation (or sometimes Liouville-von Neumann equation), and describes density matrix evolution in a closed system (identical to Schr{\"o}dinger evolution of each component)
\end{itemize}
\item Easy to tell if a system is in a pure state (i.e. can be described as a state vector), density operator has a single eigenvalue which takes value $1$ (always diagonalizable because Hermitian by construction), this can be checked for instance by taking $\mathrm{Tr}[\hat{\rho}^2]$ and seeing if it is less than $1$ (this is equivalent to asking if a system is entangled with other systems)
\item If we want to know what a subsystem within a quantum system is doing, we can take a ``partial trace'' which gives a density operator describing the subsystem, assuming a system can be divided into subsystems $A$ and $B$
\begin{equation}
\hat{\rho}_A=\mathrm{Tr}_B[\hat{\rho}]=\sum_j(1_A\otimes\bra{j}_B)\hat{\rho}(1_A\otimes\ket{j}_B)
\end{equation}
\item Nothing is done to subsystem $A$ ($1_A$ is an identity matrix), but we perform a trace using an orthonormal basis describing subsystem $B$
\item Easy to tell if two subsystems are entangled within a pure state $\ket{\psi}$, take partial trace and see if it is pure (mathematically, if $\mathrm{Tr}[\mathrm{Tr}_B[\hat{\ketbra{\psi}{\psi}}]^2]=1$), trick doesn't work if overall system cannot be described as a state vector
\item Much more to be said here, but this is probably enough to get started

\end{itemize}

\section{Measurements and noise in the density operator formalism}

\begin{itemize}
\item Since density operators are able to represent incoherent superpositions, we can consider multiple potential outcomes of a projective measurement simultaneously
\item Consider a projective measurement with $n$ unique outcomes, where we project into each outcome with a projector matrix $P_j$ which are mutually exclusive and therefore simultaneously diagonalizable
\item The effect of a projective measurement on $\hat{\rho}$ is $\hat{\rho} \rightarrow \sum^{n-1}_{j=0} P_j\hat{\rho}P^\dagger_j$
\begin{itemize}
\item The diagonal (in the diagonal basis of $P_j$) elements of $\hat{\rho}$ are unaffected, measurement cannot change probabilities, off diagonal elements between different $P_j$ will be set to zero
\item Projective measurement can be seen as setting off diagonal elements between different outcomes to zero
\item Very simple example, start with $\ketbra{+}{+}=\frac{1}{2}\left(\begin{array}{cc} 1 & 1 \\ 1 & 1 \end{array}\right)$ and measure in $\{\ket{0},\ket{1}\}$ basis, projectors will be $P_0=\left(\begin{array}{cc} 1 & 0 \\ 0 & 0 \end{array}\right)$ and $P_1=\left(\begin{array}{cc} 0 & 0 \\ 0 & 1 \end{array}\right)$
\begin{align}
P_0\ketbra{+}{+}P^\dagger_0=\frac{1}{2}\left(\begin{array}{cc} 1 & 0 \\ 0 & 0 \end{array}\right) \\
P_1\ketbra{+}{+}P^\dagger_1=\frac{1}{2}\left(\begin{array}{cc} 0 & 0 \\ 0 & 1 \end{array}\right)\\
P_0\ketbra{+}{+}P^\dagger_0+P_1\ketbra{+}{+}P^\dagger_1=\frac{1}{2}\left(\begin{array}{cc} 1 & 0 \\ 0 & 1 \end{array}\right)
\end{align}
\item Diagonal terms stay the same, but off-diagonal coherence terms are lost
\end{itemize} 
\item Can also consider continuous noise processes (the ``right way'' rather than just adding imaginary terms to Hamiltonians)
\item For continuous processes, add more terms to von Neumann equation, but we need to do this in a way which the density matrix remains a ``good'' density matrix, need to ensure two things\footnote{If you want to know the formal term for this, the operations need to be completely positive and trace preserving maps, often abbreviated CPTP}
\begin{enumerate}
\item Trace preserving, probabilities must sum to $1$
\item Positivity preserving, probabilities must always be $\ge 0$ no matter what probability is measured
\end{enumerate}
\item Preserving positivity is equivalent to guaranteeing that eigenvalues of $\hat{\rho}$ are $\ge 0$ for any valid starting state
\item These properties greatly limit the form the differential equation can take, property is only guaranteed for master equations of Lindblad form
\begin{equation}
\frac{\partial}{\partial t}\hat{\rho}=-i[\hat{H},\hat{\rho}]+\sum_{j}\gamma_{j}\left(\hat{L}_j\hat{\rho}\hat{L}^\dagger_j-\frac{1}{2}(\hat{L}^\dagger_j\hat{L}_j\hat{\rho}+\hat{\rho}\hat{L}^\dagger_j\hat{L}_j)\right) \label{Eq:Lindblad}
\end{equation}
\item Where the operators $\hat{L}_i$ describe different ways the system can interact with it's environment, and $\gamma_j$ describes the strength of the interaction
\item Simple examples include $\hat{L}_i$ consisting of single off-diagonal terms which correspond to decay from one state into another, or diagonal versions which correspond to dephasing
\item Deriving what values these terms take to model real systems is hard work, and involves a lot of approximations, I advise against doing it unless you are an ``expert'' at it (I am not an expert at this)
\item Example: spontaneous decay described by the Lindblad equation
\begin{itemize}
\item Two state atomic system, initially in $\ket{+}$ state ($\ketbra{+}{+}=\frac{1}{2}\left(\begin{array}{cc} 1 & 1 \\ 1 & 1 \end{array}\right)$), phase oscillation due to level spacing $\hat{H}=\Delta Z$, decay to lower energy state ($\ket{1}$ state) at rate $\gamma$, $\hat{L}=\left(\begin{array}{cc} 0 & 1 \\ 0 & 0 \end{array}\right)$
\item We note that $\hat{L}^\dagger \hat{L}=\left(\begin{array}{cc} 1 & 0 \\ 0 & 0 \end{array}\right)$
\item Taking $\hat{\rho}=\left(\begin{array}{cc} \rho_{00} & \rho_{01} \\ \rho_{10} & \rho_{11} \end{array}\right)$ and plugging in, we get the following set of differential equations
\begin{align}
\frac{\partial}{\partial t}\rho_{00}=-\gamma \rho_{00}\\
\frac{\partial}{\partial t}\rho_{11}=\gamma \rho_{00}\\
\frac{\partial}{\partial t}\rho_{01}=(-2 i\Delta -\gamma) \rho_{01}\\
\frac{\partial}{\partial t}\rho_{10}=(2 i\Delta -\gamma) \rho_{10}
\end{align}
\item These equations can readily be solved for the given initial conditions
\begin{equation}
\hat{\rho}(t)=\frac{1}{2}\left(\begin{array}{cc} \exp(-\gamma t) & \exp(-2 i\Delta t-\gamma t) \\ \exp(2 i\Delta t-\gamma t) & 2-\exp(-\gamma t) \end{array}\right)
\end{equation}
\item In the long time will be found in the $\ket{1}$ state with probability $1$, off diagonal elements both oscillate and decay, but decay rate depends only on bath coupling, not energy separation, $\mathrm{Tr}[\hat{\rho}(t)]=1$ at all times, remains Hermitian at all times because $\rho_{01}$ and $\rho_{10}$ are complex conjugates
\item Exercise: diagonalize and show eigenvalues remain $\ge 0$ at all times
\end{itemize}
\end{itemize}

\section{A brief digression about superoperators}

\begin{itemize}
\item Note that equation \ref{Eq:Lindblad} is linear in $\hat{\rho}$, this means that it can be written as a matrix operating on a ``vectorized'' version of $\hat{\rho}$,  $\hat{\rho}$ becomes $\vec{\rho}$ \footnote{multiple ways to do this, one way is to just take each column of the matrix and put them into a vector one after another}, the equation than becomes
\begin{equation}
\frac{\partial}{\partial t}\vec{\rho}=\mathcal{L}\vec{\rho}
\end{equation}
\item $\mathcal{L}$ is just a big matrix describing the linear transformations on the elements of the density matrix (it is also know as a ``superoperator'' if you want it to sound ``cool'' and ``scary''), this admits a formal solution
\begin{equation}
\vec{\rho}(t)=\exp(\mathcal{L} t)\vec{\rho}(0)
\end{equation}
\item Computers are generally good at exponentiating matrices, so this formal solution  is actually a good way to solve the equation numerically (for time dependent Hamiltonians this can be generalized by performing many short evolutions)
\item $\mathcal{L}$ is not Hermitian in general, so not guaranteed to be diagonalizable, but it often is diagonalizable in practice, and can always be put into a slightly more general ``almost diagonal'' form known as the Jordan form, eigenvalues (or Jordan blocks) can tell us a lot about how the system will behave
\item Quantum mechanics is linear even when open system effects are considered! (projective measurements can be written as linear operators using similar tricks, as can partial trace, as well as a ton of other operations)
\end{itemize}
\section{Zeno and adiabatic effects in the density operator formalism}
\begin{itemize}
\item Consider the differential equation for an individual matrix element under von Neumann equation
\begin{equation}
\frac{\partial}{\partial t}\hat{\rho}_{j,l}=-i\sum_{k}H_{j,k}\hat{\rho}_{k,l}+i\sum_{k}H_{k,l}\hat{\rho}_{j,k}
\end{equation}
\item Equipped with the density matrix formalism it is much easier to see the Zeno effect mathematically, consider time dependence of diagonal element under von Neumann equation
\begin{equation}
\frac{\partial}{\partial t}\hat{\rho}_{j,j}=-i\sum_{k}H_{j,k}\hat{\rho}_{k,j}+i\sum_{k}H_{k,j}\hat{\rho}_{j,k}=2\mathrm{Im}[\sum_k H_{j,k}\hat{\rho}_{k,j}]
\end{equation}
\item Where the last equality takes advantage of the fact that both $\hat{H}$ and $\hat{\rho}$ are Hermitian, so therefore $\hat{H}_{j,k}\hat{\rho}_{j,k}=\hat{H}^*_{k,j}\hat{\rho}^*_{k,j}$\footnote{Hermiticity also implies that $\mathrm{Im}[\hat{H}_{k,k}\hat{\rho}_{k,k}]=0$}
\item The derivative of diagonal elements of $\hat{\rho}$ only depends on off-diagonal elements, therefore anything which eliminates off diagonal elements will stop von Neumann equation dynamics
\item Projective measurements turn coherent superpositions to incoherent ones, kill these off diagonal elements, and stop dynamics
\item Note: from equation \ref{Eq:Lindblad}, we notice that $\hat{L}_j\hat{\rho}\hat{L}^\dagger_j$ can directly move probability around without first going through off-diagonal elements, open system operations to not need coherence to operate
\item Adiabatic theorem is also nicer to see, in the diagonal basis of the Hamiltonian, the differential equation for off diagonal elements becomes simple $\frac{\partial}{\partial t}\hat{\rho}_{j,k}=i\hat{\rho}_{j,k}(E_j-E_k)$, simple solution is only oscillating phases $\hat{\rho}_{j,k}(t)=\exp[i(E_j-E_k)t]\hat{\rho}_{j,k}(t)$
\item If phase rotation is much faster than change in basis, than interference will occur before a significant off diagonal element can be built up
\item A related trick which works under general assumptions (non-degenerate gaps, the spacing between every pair of energy levels is unique\footnote{n.b.~the harmonic oscillator does not fall under these assumptions}), is to set the off diagonal elements between all energy levels to $0$, this is effectively the average result for measuring at a random (potentially long) time, allows us to extract ``infinite time averages''
\end{itemize}




\end{document}